\begin{table}[!htbp]
\small
\setlength{\tabcolsep}{4pt}
\caption{Household inflation variance and its association with mean inflation explained by observed characteristics}
\label{tab:household-inflation-explained-characteristics}
\centering
\resizebox{\textwidth}{!}{%
\begin{tabular}{lrrrrrr}
\toprule
 & Income & Age & Education & Household size & Urbanization & All interactions \\
\midrule
Mean share of variance explained (\%) & 2.5 & 1.8 & 1.5 & 1.5 & 1.6 & 12.4 \\
Share of inflation--dispersion correlation explained (\%) & 0.4 & 0.2 & 0.1 & 0.1 & -0.2 & 0.5 \\
\bottomrule
\end{tabular}%
}
\begin{flushleft}
\footnotesize{\emph{Notes:} Separate unweighted group-mean models are estimated for each country and month from January 2021 to December 2025 using the common complete-case sample of 138,798 households in 17 countries. Household inflation is constructed at the three-digit ECOICOP1 level. Each of the first five columns classifies households only by the indicated characteristic; these columns are separate one-way partitions and are not additive. The last column uses the full interaction of income quintile, age, education, household size, and degree of urbanization. The first row averages the country-weighted national shares of variance explained across months. The second row reports the percentage reduction in the correlation between mean inflation and total variance when total variance is replaced by within-group variance. A negative value means that removing between-group differences slightly increases this correlation. Country weights are annual HICP weights normalized over the sample. The exercise is descriptive and does not assign a causal interpretation to household characteristics.}
\end{flushleft}
\end{table}
